\chapter{Testing And Quality}

The repository includes multiple verification layers because the assignment explicitly weights testing and
quality.

\section{Automated Coverage}

\begin{itemize}
  \item \textbf{Backend}: pytest coverage for auth, permissions, documents, sessions, exports, AI, and
  reset flows.
  \item \textbf{Frontend}: Vitest/React Testing Library coverage for auth pages, dashboard, management
  panels, AI panel behaviour, editor-related workflows, and realtime client behaviour.
  \item \textbf{Realtime}: dedicated tests for authenticated websocket sessions, conflict resolution,
  viewer rejection, and backend-driven revert propagation.
  \item \textbf{Integration}: a root smoke test verifies the FastAPI backend and realtime relay working
  together under shared SQLite state.
  \item \textbf{Browser}: Playwright coverage for end-to-end happy paths, plus a headless stale-AI-apply
  regression.
\end{itemize}

The root command below runs the full repository verification path:

\begin{lstlisting}[language={}]
npm run test:all
\end{lstlisting}

\section{Operational Quality}

The project also includes setup-quality features intended to keep the runtime reproducible:

\begin{itemize}
  \item root \texttt{Makefile} and \texttt{run.sh} wrappers,
  \item \texttt{.env.example} files,
  \item FastAPI OpenAPI docs,
  \item reproducible local startup,
  \item Docker Compose configuration,
  \item repository-level docs describing contracts and deviations.
\end{itemize}
